\section{Activation Functions}
\label{sec:annex-activation-functions}

This annex documents the complete activation-function library available in the
model. The library provides \textbf{43 feed-forward activation functions}: 40
elementwise activations and 3 gated-FFN variants. Activations are grouped into
five families, following the taxonomy of Dubey et al.~\citep{dubey2021activations}
and the analytic overview of Lederer~\citep{lederer2021activations}, augmented
with the learnable Rational Activation Function (RAF) of
Fang et al.~\citep{fang2023raft}.

The selection is driven by the \texttt{ffn\_activation} configuration field (a
string enum) and dispatched at runtime by a factory function that mirrors the
normalization-layer dispatch (Annex~7).

\[
\texttt{ffn\_activation} \in \{\texttt{silu},\;\texttt{gelu},\;\dots,\;
\texttt{raf},\;\texttt{swiglu},\;\texttt{geglu},\;\texttt{reglu}\}
\]

\subsection{Classical / Sigmoid--Tanh Family (12 activations)}

The classical family comprises smooth, stateless activations rooted in
logistic and algebraic functions~\citep{lederer2021activations,dubey2021activations}.
All share monotonicity or near-monotonicity and are continuously differentiable.

\begin{align}
  \text{Sigmoid}(x) &= \frac{1}{1+e^{-x}}
    & \text{Range: } (0,1),\;\text{monotonic, smooth}
    & \quad \citep{lederer2021activations} \\[4pt]
  \text{Tanh}(x) &= \frac{e^{x}-e^{-x}}{e^{x}+e^{-x}} = 2\,\sigma(2x)-1
    & \text{Range: } (-1,1),\;\text{monotonic, smooth}
    & \quad \citep{lederer2021activations} \\[4pt]
  \text{Arctan}(x) &= \arctan(x)
    & \text{Range: } (-\tfrac{\pi}{2},\tfrac{\pi}{2}),\;\text{monotonic, smooth}
    & \quad \citep{lederer2021activations} \\[4pt]
  \text{Softsign}(x) &= \frac{x}{1+|x|}
    & \text{Range: } (-1,1),\;\text{smooth, once-diff.\ at 0}
    & \quad \citep{lederer2021activations} \\[4pt]
  \text{Elliott}(x) &= \frac{x}{1+|x|}
    & \text{Alias of Softsign (survey ``elliottsig'' notation)}
    & \quad \citep{lederer2021activations} \\[4pt]
  \text{Identity}(x) &= x
    & \text{Range: } (-\infty,\infty),\;\text{linear, }C^\infty
    & \quad \text{(internal)}
\end{align}

\begin{align}
  \text{Softplus}(x) &= \log(1+e^{x})
    & \text{Range: } [0,\infty),\;\text{smooth, }C^\infty
    & \quad \citep{glorot2011softplus} \\[4pt]
  \text{Mish}(x) &= x\,\tanh\!\big(\text{softplus}(x)\big)
    = x\,\tanh\!\big(\log(1+e^{x})\big)
    & \text{Range: } [\approx{-0.31},\infty),\;\text{non-monotonic, smooth}
    & \quad \citep{misra2019mish} \\[4pt]
  \text{GELU}(x) &= x\,\Phi(x) = x\cdot\tfrac{1}{2}\Big(1+\text{erf}\!\Big(\frac{x}{\sqrt{2}}\Big)\Big)
    & \text{Range: } (\approx{-0.17},\infty),\;\text{smooth, probabilistic}
    & \quad \citep{hendrycks2016gelu} \\[4pt]
  \text{GELU-tanh}(x) &= \tfrac{x}{2}\Big(1+\tanh\!\Big(\sqrt{\tfrac{2}{\pi}}
    (x+0.044715\,x^{3})\Big)\Big)
    & \text{Tanh approximation of GELU (faster)}
    & \quad \citep{hendrycks2016gelu} \\[4pt]
  \text{ReLU}(x) &= \max(0,x)
    & \text{Range: } [0,\infty),\;\text{piecewise-linear, non-smooth at 0}
    & \quad \citep{nair2010relu} \\[4pt]
  \text{SiLU}(x) &= x\,\sigma(x) = \frac{x}{1+e^{-x}}
    & \text{Range: } (\approx{-0.278},\infty),\;\text{non-monotonic, smooth}
    & \quad \citep{elfwing2018silu}
\end{align}

\noindent\textbf{SiLU} (Sigmoid Linear Unit) is the default activation
(\texttt{ffn\_activation=silu}). When $\beta=1$, the parametric Swish reduces
to SiLU. GELU and its tanh approximation are the standard in BERT and
GPT-2/3~\citep{hendrycks2016gelu}. Mish~\citep{misra2019mish} is a
self-regularized, non-monotonic composition of Softplus and Tanh, reported to
outperform Swish on several benchmarks.

\subsection{Rectified Family (12 activations)}

The rectifier family extends ReLU with learnable slopes, bounded variants, and
rectified-hyperbolic forms~\citep{nair2010relu,dubey2021activations}.

\begin{align}
  \text{LeakyReLU}(x) &= \begin{cases} x & x\ge 0 \\ a\,x & x < 0 \end{cases},
    \quad a=0.01
    & \text{Range: } (-\infty,\infty)
    & \quad \citep{maas2013leakyrelu} \\[4pt]
  \text{ReLU6}(x) &= \min\!\big(\max(0,x),\,6\big)
    & \text{Range: } [0,6],\;\text{bounded, fixed-point friendly}
    & \quad \citep{howard2017mobile} \\[4pt]
  \text{HardSwish}(x) &= x\cdot\frac{\text{ReLU6}(x+3)}{6}
    & \text{Range: } (\approx{-1.67},\infty),\;\text{cheap Swish approx.}
    & \quad \citep{howard2019hardswish} \\[4pt]
  \text{PReLU}(x) &= \max(0,x)+\mathbf{p}\odot\min(0,x)
    & \text{Learnable per-channel slope } \mathbf{p},\;\text{init 0.25}
    & \quad \citep{he2015prelu} \\[4pt]
  \text{AbsReLU}(x) &= \max(0,x)-\max(0,-x) = x\;\text{rectified to sign-ramp}
    & \text{Range: } (-\infty,\infty)
    & \quad \citep{dubey2021activations} \\[4pt]
  \text{NLReLU}(x) &= \beta\,\log\!\big(1+\max(0,x)\big)
    & \text{Range: } [0,\infty),\;\beta=1.0,\;\text{logarithmic compression}
    & \quad \citep{dubey2021activations}
\end{align}

\begin{align}
  \text{BReLU}(x) &= \min\!\big(\max(0,x),\,t\big),\quad t=1.0
    & \text{Range: } [0,t],\;\text{bounded rectifier}
    & \quad \citep{dubey2021activations} \\[4pt]
  \text{VReLU}(x) &= |x|
    & \text{Range: } [0,\infty),\;\text{symmetric rectifier}
    & \quad \citep{dubey2021activations} \\[4pt]
  \text{Hexpo}(x) &= a\,\max(0,x)-c\,\max(0,-x),\quad a{=}c{=}1.0
    & \text{Asymmetric two-sided rectifier}
    & \quad \citep{dubey2021activations} \\[4pt]
  \text{PenalizedTanh}(x) &= \max\!\big(0,\,\tanh(x)\big)
    & \text{Range: } [0,1),\;\text{rectified tanh, smooth near 0}
    & \quad \citep{dubey2021activations} \\[4pt]
  \text{DisReLU}(x) &= \max(0,\,x-\delta),\quad \delta=0.0
    & \text{Right-shifted ReLU hinge}
    & \quad \citep{dubey2021activations} \\[4pt]
  \text{LiSHT}(x) &= x\,\tanh(x)
    & \text{Range: } [0,\infty)\;\text{(magnitude), non-monotonic, smooth}
    & \quad \citep{roy2019lisht}
\end{align}

\noindent\textbf{PReLU}~\citep{he2015prelu} makes the negative slope a
per-channel learnable parameter, enabling the network to learn the optimal
leak per feature. \textbf{HardSwish}~\citep{howard2019hardswish} is the cheap
Swish approximation used in MobileNetV3. \textbf{ReLU6}~\citep{howard2017mobile}
was introduced in MobileNetV1 for fixed-point quantization friendliness.

\subsection{Exponential / ELU Family (12 activations)}

The exponential-linear family provides smooth negative saturation and
self-normalizing variants~\citep{clevert2016elu,klambauer2017selu,dubey2021activations}.
Several members introduce per-channel learnable parameters.

\begin{align}
  \text{ELU}(x) &= \begin{cases} x & x\ge 0 \\ \alpha(e^{x}-1) & x<0 \end{cases},
    \quad \alpha=1.0
    & \text{Range: } (-\alpha,\infty)
    & \quad \citep{clevert2016elu} \\[4pt]
  \text{SELU}(x) &= \lambda\,\text{ELU}_{\alpha'}(x)
    & \alpha'{\approx}1.6733,\;\lambda{\approx}1.0507,
    & \quad \citep{klambauer2017selu} \\
  & & \text{self-normalizing (mean 0, var 1)} & \notag \\[4pt]
  \text{CELU}(x) &= \begin{cases} x & x\ge 0 \\ \alpha(e^{x/\alpha}-1) & x<0 \end{cases},
    \quad \alpha=1.0
    & \text{Once-diff.\ at 0 for any } \alpha
    & \quad \citep{barron2017celu} \\[4pt]
  \text{PELU}(x) &= \begin{cases} \tfrac{\alpha}{\beta}\,x & x\ge 0 \\
    \alpha(e^{x/\beta}-1) & x<0 \end{cases}
    & \alpha,\beta\text{ learnable per-channel}
    & \quad \citep{trottier2017pelu} \\[4pt]
  \text{MPELU}(x) &= \text{ReLU}(x)+\beta\cdot\text{ELU}_{\alpha}(x)\cdot\mathbb{1}_{x<0}
    & \alpha,\beta\text{ learnable per-channel}
    & \quad \citep{dubey2021activations}
\end{align}

\begin{align}
  \text{FELU}(x) &= \begin{cases} x & x\ge 0 \\ \alpha(e^{x}-1) & x<0 \end{cases}
    & \alpha\text{ learnable, clamped to } [0,1]
    & \quad \citep{dubey2021activations} \\[4pt]
  \text{EELU}(x) &= \beta\,\text{ReLU}(x)+\alpha(e^{x}-1)\cdot\mathbb{1}_{x<0}
    & \alpha,\beta\text{ learnable per-channel}
    & \quad \citep{dubey2021activations} \\[4pt]
  \text{PDELU}(x) &= \begin{cases} x & x\ge 0 \\
    \alpha(e^{x/\alpha}-1)+(1-\alpha)\,x & x<0 \end{cases}
    & \alpha\text{ learnable, interpolates CELU--identity}
    & \quad \citep{dubey2021activations} \\[4pt]
  \text{PREU}(x) &= \beta\,\text{ReLU}(x)+\alpha(e^{x}-1)\cdot\mathbb{1}_{x<0}^{\le 0}
    & \alpha,\beta\text{ learnable per-channel}
    & \quad \citep{dubey2021activations} \\[4pt]
  \text{SoftExp}(x) &= \begin{cases}
    \tfrac{1}{\alpha}\big(e^{\alpha x}-1\big)+\alpha & \alpha>0 \\
    x & \alpha=0 \\
    -\tfrac{1}{\alpha}\ln\!\big(1-\alpha(x+\alpha)\big) & \alpha<0
    \end{cases}
    & \alpha\text{ learnable; exp/linear/log interpolation}
    & \quad \citep{godfrey2015softexp}
\end{align}

\begin{align}
  \text{ELiSH}(x) &= \begin{cases} x\,\sigma(x) & x\ge 0 \\
    (e^{x}-1)\,\sigma(x) & x<0 \end{cases}
    & \text{Swish pos.\ branch, ELU-gated neg.\ branch}
    & \quad \citep{basirat2019elish} \\[4pt]
  \text{HardELiSH}(x) &= \begin{cases} x\cdot\text{hard\_sigmoid}(x) & x\ge 0 \\
    (e^{x}-1)\cdot\text{hard\_sigmoid}(x) & x<0 \end{cases}
    & \text{Hard-sigmoid variant of ELiSH}
    & \quad \citep{basirat2019elish}
\end{align}

\noindent where $\text{hard\_sigmoid}(x)=\text{ReLU6}(x+3)/6$.

\noindent\textbf{ELU}~\citep{clevert2016elu} provides smooth negative
saturation. \textbf{SELU}~\citep{klambauer2017selu} induces self-normalization
with fixed $\alpha'{\approx}1.6733$, $\lambda{\approx}1.0507$.
\textbf{CELU}~\citep{barron2017celu} is once-differentiable at 0 for any
$\alpha$ (unlike ELU). \textbf{PELU}~\citep{trottier2017pelu} makes both
$\alpha,\beta$ learnable per-channel. The parametric variants
(MPELU, FELU, EELU, PDELU, PREU) generalize ELU and PReLU by introducing
learnable per-channel parameters for the negative and positive branches.
\textbf{SoftExp}~\citep{godfrey2015softexp} continuously interpolates between
logarithmic ($\alpha<0$), linear ($\alpha=0$), and exponential ($\alpha>0$)
regimes. \textbf{ELiSH} and \textbf{HardELiSH}~\citep{basirat2019elish}
combine Swish and ELU behaviors via a piecewise definition.

\subsection{Swish and Parametric Activations}

Swish~\citep{ramachandran2017swish} is a family parameterized by a slope
$\beta$:
\[
  \text{Swish}_\beta(x) = x\,\sigma(\beta x) = \frac{x}{1+e^{-\beta x}}.
\]
When $\beta=1$ this is SiLU (the default). Large $\beta$ approaches
ReLU; $\beta\to 0$ approaches the linear function $x/2$.

\begin{itemize}[leftmargin=*]
  \item \textbf{Swish} (\texttt{swish}): fixed $\beta$ (default 1.0, configured
    via \texttt{ffn\_activation\_config.swish\_beta}). No learnable parameters.
  \item \textbf{SwishTrainable} (\texttt{swish\_trainable}): $\beta$ is a
    per-channel learnable parameter, initialized to 1.0. Equivalent to SiLU at
    initialization, then adapts during training.
  \item \textbf{Maxout} (\texttt{maxout}): computes $k$ linear projections and
    takes the elementwise maximum: $\max_{i=0}^{k-1}(W_i x + b_i)$.
    A universal approximator for any continuous function given enough pieces.
    Parameter count increases by factor $k$ (default $k=2$, configured via
    \texttt{ffn\_activation\_config.maxout\_pieces}).
\end{itemize}

\noindent References: Swish~\citep{ramachandran2017swish},
Maxout~\citep{goodfellow2013maxout}.

\subsection{Rational Activation Functions (RAF)}
\label{sec:raf}

Following Fang et al.~\citep{fang2023raft}, a Rational Activation Function is a
learnable ratio of two low-degree polynomials (a Pad\'e approximant):
\[
  F(x) = \frac{P(x)}{Q(x)}
        = \frac{\displaystyle\sum_{j=0}^{m} a_j\, x^{j}}
               {1 + R(x)}
\]
where $\mathbf{a}\in\mathbb{R}^{m+1}$ and $\mathbf{b}\in\mathbb{R}^{n}$ are
\emph{learnable}, and the denominator term $R(x)$ depends on the version.
The default configuration matches the paper: degree $(m,n)=(5,4)$, version
``A'' (safe), initialized by a least-squares fit to GELU on $[-3,3]$.

The implementation supports five denominator variants:

\begin{itemize}[leftmargin=*]
  \item \textbf{Version A} (default, ``safe''):
    $Q(x)=1+\sum_{k}|b_k\, x^{k+1}|$ (per-term absolute value).
    Guarantees $Q(x)\ge 1$, no division by zero.
  \item \textbf{Version B}: $Q(x)=1+\big|\sum_{k} b_k\, x^{k+1}\big|$
    (absolute value of the whole sum). Also guarantees $Q(x)\ge 1$.
  \item \textbf{Version C}: $Q(x)=0.1+\big|\sum_{k} b_k\, x^{k}\big|$
    with a $0.1$ floor. Minimum denominator $0.1$.
  \item \textbf{Version D}: like B, but injects uniform multiplicative noise
    $U(1-\varepsilon, 1+\varepsilon)$ on the denominator weights during
    \emph{training only} ($\varepsilon=0.1$). Regularization effect.
  \item \textbf{Version N} (``nunsafe''):
    $Q(x)=1+\sum_{k} b_k\, x^{k+1}$ without any absolute value.
    Can have poles (denominator zeros); use with care.
\end{itemize}

\textbf{RAFT input scaling.} Because a Pad\'e approximant is a reliable
universal approximator only on a bounded interval, the RAFT model preprocesses
each token's pre-activations with a per-token min--max scaling into $[-3,3]$:
\[
  \tilde{x} = \frac{x-\min(x)}{\max(x)-\min(x)}\cdot 6 - 3.
\]
This is toggled by \texttt{ffn\_activation\_config.raf\_input\_scaling}
(default \texttt{false}).

\textbf{Initialization targets.} The rational coefficients can be initialized by
fitting to any of the following target functions: \texttt{gelu} (default),
\texttt{relu}, \texttt{leaky\_relu}, \texttt{leaky\_relu\_0.1},
\texttt{sigmoid}, \texttt{tanh}, \texttt{swish}/\texttt{silu}, or
\texttt{identity}. The fit uses a least-squares solve on 2001 uniformly spaced
samples over $[-3,3]$.

\textbf{Freezing.} Setting \texttt{raf\_trainable=false} freezes the rational
parameters, useful for parameter-efficient fine-tuning. Frozen RAFs are
typically paired with a separate (higher) learning rate for the remaining
network parameters.

\textbf{Note on naming.} RAF stands for \emph{Rational} Activation Function,
not ``Rectified''; it is the learnable polynomial ratio described above.

\subsection{Gated FFN Variants: SwiGLU, GEGLU, ReGLU}
\label{sec:glu}

SwiGLU, GEGLU, and ReGLU~\citep{shazeer2020glu} are gated feed-forward units
that replace the standard $\text{Linear}\to\text{activation}\to\text{Linear}$
block. They own their own linear projections and are therefore implemented as
FFN modules, \emph{not} as elementwise activations:
\[
  \text{GatedFFN}(x) = \text{dropout}\!\Big(\,
    W_{\text{down}}\Big(\text{act}(x\,W_{\text{gate}})
    \odot (x\,W_{\text{up}})\Big)\Big),
\]
where $W_{\text{gate}}, W_{\text{up}}\in\mathbb{R}^{d\times d_{\text{ff}}}$ and
$W_{\text{down}}\in\mathbb{R}^{d_{\text{ff}}\times d}$. The gate function is:
\begin{itemize}[leftmargin=*]
  \item \textbf{SwiGLU} (\texttt{swiglu}): $\text{act}=\text{SiLU}$
    (the default, used in Llama, PaLM).
  \item \textbf{GEGLU} (\texttt{geglu}): $\text{act}=\text{GELU}$
    (used in some T5 variants).
  \item \textbf{ReGLU} (\texttt{reglu}): $\text{act}=\text{ReLU}$
    (cheapest option).
\end{itemize}

\noindent When \texttt{ffn\_activation} is one of
\texttt{swiglu}/\texttt{geglu}/\texttt{reglu}, the feed-forward layer swaps
the standard dense and MoE FFN blocks for a gated FFN built with the same
projection class (BitLinear under BitNet). Bias is disabled by default.

\subsection{Configuration Schema}

The \texttt{ffn\_activation\_config} object groups all learnable-activation
parameters. It is optional and ignored for stateless activations.

\begin{itemize}[leftmargin=*]
  \item \texttt{raf\_degrees}: $[m,n]$ Pad\'e numerator/denominator degrees
    (ints $\ge 1$). Default: $[5,4]$.
  \item \texttt{raf\_version}: denominator form, one of
    \texttt{A} (default), \texttt{B}, \texttt{C}, \texttt{D}, \texttt{N}.
  \item \texttt{raf\_approx\_func}: initialization target. One of
    \texttt{gelu} (default), \texttt{relu}, \texttt{leaky\_relu},
    \texttt{leaky\_relu\_0.1}, \texttt{sigmoid}, \texttt{tanh},
    \texttt{swish}, \texttt{silu}, \texttt{identity}.
  \item \texttt{raf\_trainable}: freeze the rational parameters when
    \texttt{false}. Default: \texttt{true}.
  \item \texttt{raf\_input\_scaling}: per-token min--max scaling to $[-3,3]$
    before the rational (RAFT preprocessing). Default: \texttt{false}.
  \item \texttt{prelu\_init}: initial negative slope for PReLU. Default: $0.25$.
  \item \texttt{elu\_alpha}: ELU saturation constant. Default: $1.0$.
  \item \texttt{celu\_alpha}: CELU saturation constant. Default: $1.0$.
  \item \texttt{swish\_beta}: fixed $\beta$ for Swish / initial $\beta$ for
    SwishTrainable. Default: $1.0$.
  \item \texttt{leaky\_relu\_slope}: negative slope for LeakyReLU.
    Default: $0.01$.
  \item \texttt{maxout\_pieces}: number of linear pieces $k$ for Maxout
    (int $\ge 1$). Default: $2$.
  \item Additional ELU-parametric keys:
    \texttt{pelu\_alpha}, \texttt{mpelu\_alpha}, \texttt{mpelu\_beta},
    \texttt{felu\_alpha}, \texttt{eelu\_alpha}, \texttt{eelu\_beta},
    \texttt{pdelu\_alpha}, \texttt{preu\_alpha}, \texttt{preu\_beta},
    \texttt{softexp\_alpha}. All default to $1.0$ (except
    \texttt{softexp\_alpha} which defaults to $0.0$).
\end{itemize}

\subsection{Selection Guide}

The following decision tree aids activation choice:
\begin{itemize}[leftmargin=*]
  \item \textbf{Safe default} $\to$ \texttt{silu} (SiLU/Swish$_1$, used in
    Llama, PaLM).
  \item \textbf{Classic stability} $\to$ \texttt{gelu} or \texttt{relu}.
  \item \textbf{Smooth non-monotonic} $\to$ \texttt{mish} or \texttt{swish}.
  \item \textbf{Mobile-friendly} $\to$ \texttt{hardswish} or \texttt{relu6}.
  \item \textbf{Learnable per-task shape} $\to$ \texttt{raf},
    \texttt{prelu}, \texttt{swish\_trainable}, or \texttt{maxout}.
  \item \textbf{Self-normalizing networks} $\to$ \texttt{selu} (with
    AlphaDropout).
  \item \textbf{Gated FFN (own projections)} $\to$ \texttt{swiglu},
    \texttt{geglu}, or \texttt{reglu}.
\end{itemize}

\noindent With BitNet quantization, \texttt{silu} may be more robust than
smooth activations that rely on precise gradient computation. Learnable
activations (\texttt{raf}, \texttt{prelu}, \texttt{swish\_trainable}) add
parameters but can adapt to the task distribution.